% =====================================================================
%  Lab 3 -- Transforming User Stories to Requirements
%  ZoneSense: Smart Zone Environmental Monitoring and HVAC Control System
% =====================================================================
\documentclass[11pt,a4paper]{article}
% =====================================================================
%  ZoneSense lab reports -- shared preamble
%  CS3151 Software Engineering, Effat University, Fall 2026
%  Compiler: pdfLaTeX (Overleaf default)
% =====================================================================

% ---------- Fonts and encoding ----------
\usepackage[T1]{fontenc}
\usepackage[utf8]{inputenc}
\usepackage{lmodern}               % Latin Modern for typewriter text
\usepackage{mathptmx}              % Times for body text
\usepackage[scaled=0.90]{helvet}   % Helvetica for table headers and captions
\usepackage{textcomp}
\usepackage{amssymb}
\usepackage{microtype}

% ---------- Page layout ----------
\usepackage[a4paper,top=2.2cm,bottom=2.2cm,left=2.2cm,right=2.2cm,headheight=16pt]{geometry}
\setlength{\parindent}{0pt}
\setlength{\parskip}{6pt plus 1pt}
\linespread{1.08}
\setlength{\emergencystretch}{2em}

% ---------- Colours (taken from the SRS) ----------
\usepackage[table]{xcolor}
\definecolor{znavy}{HTML}{173A5E}   % headings, table headers
\definecolor{zteal}{HTML}{1A9EAD}   % sub-headings
\definecolor{zrow}{HTML}{F0F5FA}    % alternating table rows
\definecolor{zline}{HTML}{B7C6D5}   % table rules
\definecolor{znote}{HTML}{FFF6DB}   % call-out background
\definecolor{znotebar}{HTML}{D9A21B}% call-out bar
\definecolor{zpink}{HTML}{FBEDED}
\definecolor{zmint}{HTML}{E3F4F1}
\definecolor{zgreen}{HTML}{1E7B3A}
\definecolor{zcode}{HTML}{F5F7F9}

% ---------- Packages ----------
\usepackage{graphicx}
\usepackage{array}
\usepackage{tabularx}
\usepackage{xltabular}
\usepackage{makecell}
\usepackage{enumitem}
\usepackage{titlesec}
\usepackage{fancyhdr}
\usepackage{float}
\usepackage{caption}
\usepackage{needspace}
\usepackage[most]{tcolorbox}
\usepackage{listings}
\usepackage{xurl}
\usepackage[hidelinks]{hyperref}
\urlstyle{same}

\graphicspath{{figures/}}

% ---------- Headings ----------
\setcounter{secnumdepth}{3}
\titleformat{\section}{\LARGE\bfseries\color{znavy}}{\thesection.}{0.5em}{}
\titleformat{\subsection}{\Large\bfseries\color{zteal}}{\thesubsection}{0.6em}{}
\titleformat{\subsubsection}{\large\bfseries\color{znavy}}{\thesubsubsection}{0.6em}{}
\titleformat{\paragraph}[runin]{\normalsize\bfseries\color{znavy}}{}{0pt}{}
\titlespacing*{\section}{0pt}{6pt}{12pt}
\titlespacing*{\subsection}{0pt}{14pt}{6pt}
\titlespacing*{\subsubsection}{0pt}{12pt}{4pt}
\titlespacing*{\paragraph}{0pt}{8pt}{0.6em}

% ---------- Lists ----------
\setlist{itemsep=2pt,topsep=3pt,parsep=0pt}

% ---------- Tables ----------
\newcolumntype{L}[1]{>{\raggedright\arraybackslash}p{#1}}
\newcolumntype{C}[1]{>{\centering\arraybackslash}p{#1}}
\newcolumntype{Y}{>{\raggedright\arraybackslash}X}
\newcolumntype{B}[1]{>{\raggedright\arraybackslash\leavevmode\bfseries\color{znavy}}p{#1}} % bold navy first column
\newcolumntype{Z}{>{\centering\arraybackslash}X}
\renewcommand{\arraystretch}{1.25}
\setlength{\tabcolsep}{5pt}
\arrayrulecolor{zline}
\newcommand{\hd}[1]{\leavevmode\textcolor{white}{\sffamily\bfseries #1}}   % header cell
\newcommand{\headrow}{\rowcolor{znavy}}                            % header row
\newcommand{\stripe}{\rowcolors{2}{white}{zrow}}                    % zebra rows

% ---------- Captions ----------
\captionsetup{font={small,sf},labelfont={bf,color=znavy},textfont={bf,color=znavy},
  labelsep=endash,justification=raggedright,singlelinecheck=false,skip=4pt}
\captionsetup[figure]{font={small,sf},labelfont={color=gray},textfont={color=gray},
  justification=centering,singlelinecheck=true,skip=6pt}

% ---------- Handy macros ----------
\newcommand{\COtwo}{CO\textsubscript{2}}
\newcommand{\tick}{\ensuremath{\checkmark}}
\newcommand{\gtick}{\textcolor{zgreen}{\ensuremath{\checkmark}}}
\newcommand{\cbox}{\ensuremath{\square}}

% Call-out box (yellow, left bar) -- used for notes in the SRS
\newtcolorbox{zcallout}{enhanced,breakable,colback=znote,colframe=znote,
  borderline west={3pt}{0pt}{znotebar},arc=0pt,outer arc=0pt,
  left=8pt,right=8pt,top=5pt,bottom=5pt,fontupper=\small,before skip=8pt,after skip=8pt}

% Light-blue information box
\newtcolorbox{zinfo}{enhanced,colback=zrow,colframe=zrow,arc=0pt,outer arc=0pt,
  left=8pt,right=8pt,top=6pt,bottom=6pt,fontupper=\small,before skip=8pt,after skip=8pt}

% Code listings
\lstset{basicstyle=\ttfamily\scriptsize,columns=fullflexible,keepspaces=true,
  breaklines=true,breakatwhitespace=false,showstringspaces=false,upquote=true,
  frame=single,rulecolor=\color{zline},backgroundcolor=\color{zcode},
  xleftmargin=4pt,xrightmargin=4pt,framexleftmargin=4pt,framexrightmargin=4pt,
  aboveskip=8pt,belowskip=8pt,tabsize=4,
  literate={≤}{{$\le$}}1 {≥}{{$\ge$}}1 {→}{{$\rightarrow$}}1 {–}{{--}}1}

% ---------- Header and footer ----------
\pagestyle{fancy}
\fancyhf{}
\fancyhead[L]{\small\sffamily\textcolor{znavy}{\ShortTitle}}
\fancyhead[R]{\small\sffamily\textcolor{gray}{CS3151 -- Software Engineering}}
\fancyfoot[L]{\footnotesize\sffamily\textcolor{gray}{CS3151/2026 \quad Dr.\ Passent Elkafrawy}}
\fancyfoot[R]{\footnotesize\sffamily\textcolor{gray}{\thepage}}
\renewcommand{\headrulewidth}{0.4pt}
\renewcommand{\headrule}{\hbox to\headwidth{\color{zline}\leaders\hrule height \headrulewidth\hfill}}

% ---------- Cover page ----------
% Set \LabTitle, \ProjectTitle and \TeamMembers in main.tex before \begin{document}.
\newcommand{\makecover}{%
  \begin{titlepage}
    \centering
    \vspace*{1.2cm}
    \includegraphics[width=3.6cm]{effat_logo.png}\par
    \vspace{1.6cm}
    {\Large\bfseries Course: CS3151 -- Software Engineering\par}
    \vspace{0.35cm}
    {\large\bfseries Project -- Fall 2026\par}
    \vspace{0.15cm}
    {\large\bfseries Instructor: Dr.\ Passent Elkafrawy\par}
    \vspace{1.6cm}
    {\large\bfseries \ProjectTitle\par}
    \vspace{0.9cm}
    {\color{znavy}\rule{0.55\linewidth}{0.6pt}\par}
    \vspace{0.45cm}
    {\LARGE\bfseries\color{znavy} \LabTitle\par}
    \vspace{0.35cm}
    {\color{znavy}\rule{0.55\linewidth}{0.6pt}\par}
    \vfill
    {\large\bfseries Done By:\par}
    \vspace{0.4cm}
    {\large \TeamMembers\par}
    \vspace{2.2cm}
  \end{titlepage}%
}

% Keep a run-in heading together with the text that follows it
\newcommand{\keephead}[1]{\needspace{6\baselineskip}\paragraph{#1}\leavevmode\par\vspace{-\parskip}}

% Never leave a heading stranded at the bottom of a page
\usepackage{etoolbox}
\pretocmd{\section}{\needspace{8\baselineskip}}{}{}
\pretocmd{\subsection}{\needspace{6\baselineskip}}{}{}
\pretocmd{\subsubsection}{\needspace{5\baselineskip}}{}{}


\newcommand{\ShortTitle}{ZoneSense -- Lab 3: Requirements Specification}
\newcommand{\LabTitle}{Lab 3: Transforming User Stories to Requirements}
\newcommand{\ProjectTitle}{ZoneSense: Smart Zone Environmental Monitoring and HVAC Control System}
\newcommand{\TeamMembers}{%
  Remas Rawas --- ID: S24109161\\[6pt]
  Bayan Abduldaym --- ID: S24109196\\[6pt]
  Ghala Bin Salim --- ID: S24109159}

\counterwithin{table}{section}

% Requirement table: ID | Requirement
\newcommand{\reqhead}[1]{%
  \caption{#1}\\
  \hline\headrow \hd{ID} & \hd{Requirement}\\\hline
  \endfirsthead
  \hline\headrow \hd{ID} & \hd{Requirement}\\\hline
  \endhead}
% Non-functional requirement: name in bold, statement below
\newcommand{\nfr}[2]{\textbf{#1}\newline #2}

\begin{document}
\makecover

% ---------------------------------------------------------------------
\section{Requirements Analysis}

\subsection{Functional Requirements}

{\small\stripe
\begin{xltabular}{\linewidth}{|B{1.3cm}|Y|}
\reqhead{Functional requirements}\label{tab:fr}
FR1 & The system shall allow an authorized Building Resident to configure and uniquely identify up to 20 monitoring zones within one building, with each zone assigned a unique ID and name.\\\hline
FR2 & The system shall receive 4 environmental measurement types---temperature, \COtwo, humidity, and sound level---from each configured zone at least once every 60 seconds during normal sensor operation.\\\hline
FR3 & The system shall store 100\% of received sensor readings with the associated zone ID, sensor type, measured value, and timestamp.\\\hline
FR4 & The system shall display the most recently received temperature, \COtwo, humidity, and sound-level readings for each monitored zone within 8 seconds of receiving updated data.\\\hline
FR5 & For each newly received environmental reading, the system shall compare the value with its configured thresholds and determine its measurement severity, and recompute the zone status using FR6 within 5 seconds.\\\hline
FR6 & The system shall display exactly one zone status using this priority: critical if any available sensor is critical; otherwise warning if any available sensor is warning; otherwise unavailable if any configured sensor is unavailable; otherwise normal. The display shall update within 8 seconds of a status change. Each sensor's availability shall be shown separately, so missing data never hides a known critical or warning condition.\\\hline
FR7 & The system shall create an alert-history record for every detected threshold violation within 45 seconds. It shall display a new alert notification within 45 seconds for the first violation, a severity change, or a continuing violation after the repeat interval expires. Unchanged-severity repeat notifications within the interval shall be suppressed under FR22; their history records shall still be stored.\\\hline
FR8 & 100\% of generated environmental alerts shall contain the affected zone, abnormal environmental measurement, current reading, and alert severity when displayed.\\\hline
FR9 & For each detected abnormal temperature, \COtwo, or humidity condition that can be addressed through a supported control function, the system shall determine and display one proposed control action within 10 seconds of identifying the controllable condition.\\\hline
FR10 & The system shall require authorized Building Resident approval for 100\% of proposed control actions and shall execute 0 proposed control actions before approval is recorded.\\\hline
FR11 & The system shall allow the authorized Building Resident to approve or reject 100\% of pending proposed control actions, and shall record the selected decision within 2 seconds of submission.\\\hline
FR12 & When the authorized Building Resident approves a proposed control action, the system shall send the approved action to the supported control function within 2 seconds of approval.\\\hline
FR13 & When the authorized Building Resident rejects a proposed control action, the system shall execute 0\% of rejected actions and mark the proposal as rejected within 2 seconds of the decision.\\\hline
FR14 & The system shall record 100\% of Building Resident approval and rejection decisions with the associated proposed control action, Building Resident identifier, decision, and timestamp within 5 seconds of the decision.\\\hline
FR15 & After an approved control action or recorded maintenance activity, the system shall continue receiving readings from the affected zone at intervals of no more than 60 seconds and update its environmental status within 5 seconds of each new reading.\\\hline
FR16 & When an HVAC Maintenance Technician opens an environmental incident, the system shall display the affected zone, abnormal condition, relevant environmental readings, alert details, and associated control-action information within 10 seconds.\\\hline
FR17 & The system shall allow the HVAC Maintenance Technician to record the maintenance work performed, its status, and whether the technical problem was resolved, and shall store 100\% of submitted maintenance records within 5 seconds of submission.\\\hline
FR18 & The system shall store 100\% of generated environmental alerts, Building Resident approval and rejection decisions, executed control actions, and recorded maintenance activities for at least 12 months.\\\hline
FR19 & The system shall allow the Building Resident to retrieve and view previous alerts, approval decisions, and executed control actions associated with a selected zone within 10 seconds of the request.\\\hline
FR20 & The system shall require successful authentication before granting access to 100\% of protected system functions and shall grant or deny an authentication attempt within 5 seconds of valid credential submission.\\\hline
FR21 & For 100\% of protected function requests, the system shall verify the authenticated user's assigned role before granting access and shall deny 100\% of requests for functions not permitted to that role.\\\hline
FR22 & The authorized Building Resident shall be able to set a repeat-notification interval for each zone and measurement type from 1 to 60 minutes (default 10 minutes), measured from the last displayed notification. During a continuing abnormal condition of unchanged severity, the system shall suppress repeat notifications within this interval, update the zone status within 8 seconds of each reading, and record every violation in alert history with its notification disposition. A severity change shall bypass suppression and be notified within 45 seconds. A normal reading ends the abnormal episode; its next violation shall be treated as a first violation.\\\hline
\end{xltabular}}

\subsection{Non-Functional Requirements}

{\small\stripe
\begin{xltabular}{\linewidth}{|B{1.3cm}|Y|}
\reqhead{Non-functional requirements}\label{tab:nfr}
NFR1 & \nfr{Performance}{The system shall update the monitoring interface with newly received environmental sensor data within 8 seconds of receiving the data during normal operation.}\\\hline
NFR2 & \nfr{Alert Response Time}{The system shall record 100\% of detected threshold violations in alert history within 45 seconds. It shall display 100\% of notifications eligible under FR7 and FR22 within 45 seconds of detection; suppressed repeat notifications are exempt from display, but not from history recording.}\\\hline
NFR3 & \nfr{Availability}{If a configured sensor fails to provide data for 3 consecutive 60-second reporting intervals, the system shall mark that sensor unavailable within 8 seconds after the third missed interval, show its availability separately, and recompute the zone status using FR6. Stale readings shall not be used as current readings; receipt of fresh data shall restore that sensor's availability and trigger status recomputation.}\\\hline
NFR4 & \nfr{Security}{The system shall require successful authentication for 100\% of attempts to access protected monitoring, maintenance, approval, configuration, and control functions.}\\\hline
NFR5 & \nfr{Authorization}{The system shall verify the authenticated user's assigned role for 100\% of protected function requests and shall permit 0 requests for functions not authorized for that role.}\\\hline
NFR6 & \nfr{Auditability}{The system shall create an audit record for 100\% of approval decisions, executed control actions, and maintenance activities, containing the responsible user, timestamp, action or decision, and result, within 5 seconds of the event.}\\\hline
NFR7 & \nfr{Data Integrity}{The system shall preserve the correct zone ID, sensor type, measured value, and timestamp for 100\% of stored sensor readings, with 0 stored readings permitted to exist without all four required attributes.}\\\hline
NFR8 & \nfr{Reliability}{A failure or loss of data from any single sensor shall not interrupt the processing or display of data from the remaining functioning sensors and zones, and functioning sensor data shall continue to be displayed within the normal 8-second interface-update target.}\\\hline
NFR9 & \nfr{Usability}{The monitoring interface shall use the same four zone status categories---normal, warning, critical, and unavailable---and the FR6 priority on every screen displaying zone status. Each configured sensor's availability shall be visible separately on the zone-detail view.}\\\hline
NFR10 & \nfr{Privacy}{The sound-monitoring component shall process and store sound-level values only, and the system shall record or store 0 audio recordings and 0 audio conversations during operation.}\\\hline
NFR11 & \nfr{Maintainability}{An authorized Building Resident shall be able to change environmental thresholds and monitoring-zone configuration through the system interface without modifying source code, and a saved configuration change shall become active within 5 seconds.}\\\hline
NFR12 & \nfr{Scalability}{The system shall support up to 20 monitoring zones with 4 environmental measurement types per zone, for up to 80 environmental measurement streams, without redesign of the core monitoring architecture and while maintaining the specified interface and alert performance targets.}\\\hline
NFR13 & \nfr{Transport Security}{The system shall encrypt 100\% of connections between the browser and the application server, the sensor gateway and the application server, and the application server and the database using TLS 1.2 or higher, and shall store 0 user passwords in plain text (hashed storage only), in the prototype deployment.}\\\hline
\end{xltabular}}

% ---------------------------------------------------------------------
\clearpage
\section{Requirements Elicitation Methods}
The requirements for the proposed ZoneSense system were developed using the following approaches:

\subsection{Competitor Analysis}
Existing building management solutions from Siemens, Johnson Controls, Schneider Electric, and Honeywell were reviewed to identify common monitoring, alerting, and HVAC control capabilities.

\subsection{Review Analysis}
Publicly available user feedback and reviews of building management solutions were examined to identify potential usability concerns and operational challenges.

\subsection{Persona and Scenario Analysis}
Two personas were developed: the Building Resident and the HVAC Maintenance Technician. Their responsibilities, goals, and challenges were used to create usage scenarios and user stories.

\subsection{User Story Analysis}
The user stories were translated into functional and non-functional requirements and refined using SMART criteria to improve clarity, measurability, and testability.

\subsection*{Planned Validation Activities}
The following activities are proposed for later project stages and have not yet been conducted:
\begin{itemize}
\item \textbf{Stakeholder Interviews:} Gather feedback from facility management and HVAC maintenance personnel.
\item \textbf{Prototype Evaluation:} Present a working prototype to potential users to assess the proposed workflow and identify necessary improvements.
\end{itemize}
These activities will help validate the requirements against actual stakeholder needs.

\end{document}
